\documentclass[a4paper, 12pt]{article}

\usepackage[utf8]{inputenc}
\usepackage[T1]{fontenc}
\usepackage{textcomp}
\usepackage{amssymb}
\usepackage{newtxtext} \usepackage{newtxmath}
\usepackage{amsmath, amssymb}
\newtheorem{problem}{Problem}
\newtheorem{example}{Example}
\newtheorem{lemma}{Lemma}
\newtheorem{theorem}{Theorem}
\newtheorem{problem}{Problem}
\newtheorem{example}{Example} \newtheorem{definition}{Definition}
\newtheorem{lemma}{Lemma}
\newtheorem{theorem}{Theorem}
\DeclareMathAlphabet{\mathcal}{OMS}{cmsy}{m}{n}



\begin{document}

\textit{El sueño} 

~ 


Bajo las tenues y pequeñas luces que ofendían al aire nocturno del patio, como
si fuera un signo ilegible o un enigma, su rostro guaranítico parecía invadido
de fantasmas y vestigios. El piso de ladrillo era cercado por estrechas franjas
de tierra fértil, aledañas a los muros y alfombradas de perladas piedrecitas
blancas, desde las cuales decenas de plantas tropicales conjuraban una especie
secreta de encanto similar al de sus ojos negros. En efecto, había en aquellos
ojos todo lo que en América creemos —o fingimos creer— haber enterrado,
otorgándole lo que alguien describió correctamente como una «mirada de cinco
siglos». Sentada en silencio, descalza y con un cigarrillo en la mano, la
salvaje hojarasca de un jazmín brasilero abierta a poca distancia de su oscuro
cabello, pensé que ella también se preguntaba quiénes fueron los hombres, los
infinitos hombres, que la habían llevado allí. Ella no lo sospechaba, pero éramos
—como todas las cosas— apenas algo más que un eco. Yo, por otro lado, que era
tímido e infeliz, la contemplaba con el silencio respetuoso de los que visitan
un cementerio, disimulando mi distante asombro, aunque todo en mí sintiera que
un coro de voces antiguas afloraría, en cualquier momento, desde sus labios.

~

Por aquel entonces, todavía era un niño enamoradizo y me entregaba —como a
tantas otras cosas que la gente gustaba de llamar «pérdidas de tiempo»— a la
lectura de Matsuo Bashō y Kobayashi Issa. Parece mentira, pero aquella tarde en
que la conocí había recordado un dulce haiku que promete la incipiencia del
amor:

~
\begin{quote}
Under the cherry blossoms\newline
strangers are not\newline
really strangers.
\end{quote}

~

Diez años después, poco después de que yo terminara mi doctorado en Estados
Unidos, Paulina se entregó a la muerte. No fui a su velorio y, por varios años,
no visité su tumba. Mi vida continuó con total normalidad y podría decirse
incluso que fui feliz. Pero, como debía suceder, no pasó mucho tiempo hasta que
una mujer, bastante mayor que yo, me sugirió su rostro. Tenía una belleza
imperfecta, desordenada, y sus ojos exepresaban una voluntad invencible. La
cortejé sabiendo que cortejaba a Paulina; busqué a Paulina y me convencí de
haberla encontrado, ignorando deliberadamente la locura de esta idea. Seis meses
después esta mujer me abandonó, y, curiosamente, fue su abandono lo que
materializó, por así decirlo, el abandono de Paulina. Es curioso que en la vida
no asumamos un evento sino hasta confrontarlo en un símbolo o un símil. Hasta
ese entonces, fingía que Paulina estaba en algún lugar del mundo, destruyéndose
tal vez, como siempre lo hizo; consumiéndose lentamente, melancólica y febril,
pero viva. Pero ahora sentía su muerte y, sobre todo, sentía un reclamo, un
llamado que no podía ignorar. Era una voz que no era del todo la de Paulina,
que sólo puedo describir como la voz de su entierro y de la hierba vengativa,
que alteraba la de Paulina del mismo modo que el eco altera la voz de un hombre
que conversa solo. Mi vuelo a la Argentina despegó una semana después.

~

Fui recibido en mi patria con inesperada alegría. Aunque el calor era
insoportable y el día era clarísimo, quienes conocíamos el inclemente clima de
aquella región del mundo sabíamos que se avecinaba una tormenta tropical.
Estábamos reunidos en un antiguo patio donde la gente de la ciudad se amontona a
escuchar música bajo la amorosa sombra de los mangos y los urunday, meciendo
nuestros sentimientos al son de unas canciones con un dejo folklórico y feliz,
como se mece a un niño febril que quiere llevarse a las serenas regiones del
sueño. Pero lentamente los tintes carmesíes del arrebol fueron invadidos por una
oscuridad de plata. Los cuerpos, los cuerpos de las gentes y las cosas,
aparecían ante mí como fantasmas nativos de un mundo mágico y perfecto. Una
profunda soledad invadió mi corazón, tal vez porque la música se había vuelto
melancólica, tal vez porque en mis puños cerrados todavía me aferraba a un
puñado de ceniza y la sentía escaparse de mis manos lentamente, como el polvo
fino de un reloj de arena pasa de una parte a otra. «Así me escondo en este
mundo»—pensé—«con esta arena, con este recuerdo apretado entre mis manos».

~

Salí a sentarme solo en un banco de madera, cerca de un árbol de naranjos, con
el alma oprimida por la añoranza de una mujer que ha muerto, del deseo de verla
bailar en la hojarasca con aquellas personas que no la conocieron, del ansia de
volver a sentir sus manos y su voz llamándome, pero esta vez sin lágrimas ni
pena. ¿Qué puede doler más que el pensamiento desesperado de querer hundir los
labios en la tierra para llamar un nombre que, día a día, es recordado por menos
y menos personas? Esto pensaba, romantizando mi melancolía, mientras la tormenta
iba anunciándose en la voz del viento.

~

Sentí aproximarse, más que un cuerpo, una mirada. No sé cómo, pero aquellos ojos
supieron recordarme que había otro mundo dentro, un mundo de personas vivas que
bailaban vivamente, besándose los labios o tocándose las manos, amándose en
secreto aunque fuera por sólo unos instantes, bajo el influjo de una música a la
vez citadina y de interior. Aquellos ojos negros eran ojos vivos y, en cierto
modo, me llamaron al mundo de los vivos. Y entonces abandoné mi pensamiento,
abandoné el recuerdo de mi amada muerta, y mirando a la muchacha que ya se
levantaba de mi lado pensé:

\begin{quote}
Bajo el árbol de naranjos\newline
los extraños no son\newline
verdaderamente extraños
\end{quote}

Poco después la tormenta desató su furia, como una bendición, sobre nosotros.

~ 

Ligeramente ebrio, todavía sin querer recordarme qué me había llevado de vuelta
a mi ciudad natal, aquella noche dormí en la hamaca paraguaya del patio de la
casa de mi padre, que me condujo a un universo febril e intoxicante. Como a
través de una cámara, me vi de pie frente al río, castigado por un celaje
arrebolado que casi auguraba una especie rara y tropical de apocalipsis. Miré el
agua y sentí: «este fue el primer elemento». La voz de Paulina habló desde el
fondo del río: «aquí se fue \textit{che rekové}». (Sentí en mi corazón tanta
tristeza.) De repente, el agua dijo: «\textit{Rohechaga'u, che mborayhumí: ani
ñe ñaña y vení a zambullirte conmigo jeýna}». Me hundí en el agua lentamente, y
cuanto más entraba en la negrura de la muerte más sentía su amor, como un haz de
luz suspendido al otro lado de la vida. Bajo el agua, una voz decía una y otra
vez: «\textit{rojhayjhueteiva}». Apreté los puños y los dientes. Entonces
desperté. Y al despertar quise morirme.


~


Las frases en guaraní combinan locuciones de mi ciudad natal y lugares comunes
que se escuchan en guaranias y canciones populares. Pero la esencia de la
cuestión ---aunque no lo viera así en aquel momento--- yacía en aquello que fue
al mismo tiempo el último sentimiento del sueño y el primero de la vigilia: que
al despertar quise morirme. (Sentimos que el sueño es una forma de la muerte y
que la vida es la vigilia: yo intuí, aunque todavía veladamente, precisamente lo
contrario.) 

~

Desde aquel sueño, mi vida lentamente se transformó en la de un perro que
aguarda, serenamente triste, su periódico paseo por un jardín que, sin duda, le
parece ilimitado. Perdí todo interés en visitar su tumba e incluso en considerar
su muerte. Cada instante de mi vida consciente, lentamente, se transformó en una
espera. Deseaba soñar, porque en mis sueños o bien encontraba a Paulina o bien
podía buscarla sin la engañosa convicción de su muerte, que la vigilia no puede
suspender. Vivía mecánica y apáticamente: el mundo me era indiferente, la ciudad
donde vivía me era ajena, pero soñaba ricos universos y conocía bien hasta
sus últimos rincones, porque en ellos buscaba, y hasta a veces encontraba, a
Paulina. 

~

Muchos años antes, cuando abracé, con enferma obstinación, el estudio
científico; cuando viajé a Pensilvania para doctorarme; cuando, sobre todo, me
especialicé en el estudio del sueño, no sospeché jamás el vínculo secreto que
ataba cada uno de estos actos a su muerte, que ni siquiera había ocurrido. Pero
en cuanto aquel primer sueño me dio una intuición, todavía apenas barruntada, de
lo que debía hacer, sospeché un hilo secreto que unía la rígida estructura de mi
vida científica con una pasión secreta, estúpida e irracional, que
inexplicablemente precedió a su propia causa. Esto puede parecer extraño, pero
en verdad es algo usual en la vida de todas las personas: cada acontecimiento
determina no sólo el presente y el futuro, sino que opera sobre
el pasado y lo transforma. Descubrí, al despertar, que todas mis decisiones
anteriores a la muerte de Paulina tenían precisamente aquella muerte como
objeto, y concebí en el sitio más recóndito de mi alma que mi vida, desde aquel
instante, tenía solo una dirección posible.

~

La primera parte del experimento, curiosamente, no se distinguió demasiado de
ciertos experimentos usuales en la neurociencia del sueño, y el lector sabrá
disculpar que dedique unas pocas palabras a explicarlos. Estos experimentos
consisten en enseñar a sujetos de interés cómo soñar lúcidamente, cómo preservar
su voluntad en los sueños. Una vez la habilidad es adquirida, se les pide que
produzcan cierta acción física siempre que aparezca un estímulo de cierta clase.
La acción, por ejemplo pestañear tres veces, es visible en el
electroencefalograma y permite de este modo una especie de comunicación entre el
sujeto durmiente y los investigadores. 

~

Afortunadamente, en mi caso personal el periodo de aprendizaje fue rápido y, al
cabo de poco tiempo, podía soñar lúcidamente con poco esfuerzo. Durante un año,
registré la actividad nerviosa de mis sueños; en cada noche busqué o fabriqué a
Paulina, y registré con tres pestañeos su aparición ante mí. Eventualmente
determiné, o creí determinar, la signatura nerviosa, aunque indirecta, que
correspondía a ella: sus ojos, los sutiles rizos de su cabello rubio, la curiosa
forma de hermosura afable y mansa que caracterizaba su fuero externo, la
indomable irregularidad de su corazón. Estas manifestaciones oníricas pero, en
última instancia, perceptibles, se correspondían con una parametrización
específica de la actividad de mis neuronas, o al menos tal era mi hipótesis: la
contribución de distintas frecuencias a las oscilaciones del
electroencefalograma, la sincronización espaciotemporal de distintas regiones
del cerebro, millones de neuronas cuya actividad organizada resucitaban a una
mujer que ha muerto.

~

Pero el cerebro es un sistema caótico: un pequeño cambio en su configuración
produce una diferencia extraordinaria en su comportamiento, y no lograba
construir precisamente el sueño que buscaba. Es cierto que encontraba a Paulina,
pero no siempre era del todo Paulina, o no siempre era la Paulina que yo quería
ver: ciertas veces la conjuraba pero sus ojos eran grises, taciturnos; otras era
verdaderamente ella, pero en una tormenta o en un universo de sombras; otras
estaba allí, más allá de una frontera que ni ella ni yo podíamos cruzar; la
mayor parte de las veces era como un cuerpo sin alma, o como una imagen sin
cuerpo, o como una idea sin imagen alguna que le corresponda. 

~

Más aún, hablaba con estas apariciones y cada cual recordaba una vida distinta.
Una de ellas, que me esperaba en el río, se lamentaba de haber muerto sola en un
país desconocido, mientras que Paulina murió en su casa familiar; otra, que
existía en un mundo primitivo y salvaje, juraba nunca haber sido mi novia; otra,
que caminaba circularmente en la arena, decía que nunca había muerto. Trataba de
explicarles que no, que nos habíamos amado, que había muerto en su patio y por
obra de ella misma, y me miraban con total indiferencia, como si mis palabras no
tuvieran peso alguna en esa dimensión desconocida. Es como si, en mi intento de
resuscitarla, sólo lograra explorar un espacio infinito de Paulinas posibles,
ninguna de las cuales era la que yo había amado.

~

Noté además que cuanto menos similar era el universo del sueño al mundo en el
que Paulina y yo vivimos, más probable era que la misma Paulina no fuera del
todo ella, como si la recreación perfecta de su cuerpo y de su vida dependiera
de la de nuestra ciudad natal, de la luz estival que solía alumbrarnos, del río
y las arenas mustias que ensuciaban nuestros pies. Comprendí que no bastaba
conjurar a Paulina: era necesario conjurar el río, y no sólo el río sino el
preciso tono carmesí que el celaje tropical dotaba al río cuando escapábamos de
la escuela, y el canto específico de las aves de nuestra región del mundo, y el
húmedo calor tropical que agobiaba nuestros cuerpos. Y lentamente aprendí a
conjurar todas estas cosas, hasta obtener no solo un prototipo de Paulina sino
de una perfecta imagen del pasado, de \textit{todo} un pasado, que supimos
compartir.

~



~

~ 


\pagebreak 


no sospeché que el velo entre la
vida y la muerte se parece al que separa al sueño y la vigilia en un sentido
mucho más que literario; no sospeché, en fin, que incluso el más insignificante
de mis actos conducía al prodigio de su resurrección.


~



Este fue mi último sueño. Puedo sentirlo, pero no sé por qué. Nuestro propio
corazón es insondable: ¿tal vez he alcanzado la paz o me aproximo al olvido?
Siento que cada uno de estos sueños era el agua de un pozo tempestuoso y que
está pronto a secarse. El alma del hombre es un aljibe donde brotan las algas de
la vida. Orígenes intituló una homilía sobre el Génesis: \textit{Nuestra alma,
pozo de agua viva}. Sí... Mi pensamiento de Paulina es como un loto flotando en
el recuerdo. Ya no puedo conducirlo. Carece de todo vicio y toda culpa. Ya he
atendido sus flores y su estanque. Y es mi deber abandonarlo.

























\end{document}



